% /root/monodromy, 19 November 2026. Complete candidate for independent audit.
\subsection{Degree two on every fixed nonordinary principal surface}
\label{subsec:nonordinary-ppav-degree-two-wreath-theta-candidate}
\status{Candidate; complete proofs written for independent checks.}

Until the final theorem let $k$ be algebraically closed of characteristic
two, let $(A,\Theta)$ be an indecomposable principally polarized abelian
surface, and put $f=f(A)\in\{0,1\}$ and $L=\mathcal O_A(3\Theta)$.
We may take $\Theta$ symmetric. A different representative of the
principal polarization is a translate and gives an isomorphic family.
The known rank and covector theorems give, in degree two,
\[
 \dim V=33+f,\qquad \dim(V^*)^{\pi_1}=1+f,
 \qquad \dim K=32,
\]
where $K$ is the common kernel of all invariant covectors. We use the
already proved full degree-one group $\GL_8(\mathbf F_2)$ on every
fixed nonordinary principal surface and every Picard twist of $L$.
For a degree-one curve, ``sharp'' means that its two-rank is $8+f$.

\begin{lemma}[Lifted intersections on a fixed cubic-theta curve]
\label{lem:nonordinary-ppav-two-lifted-intersections-candidate}
Fix $\alpha\in\Pic^0(A)$ and a smooth sharp curve
$C\in|L\otimes\alpha|$. If $\widetilde C\to C$ is a connected
finite etale torsor with finite elementary abelian two-group $G$,
then the monodromy of its fibres over the eighteen points $C\cap D$,
as $D$ varies through smooth sharp transverse members of
$|L\otimes\alpha^{-1}|$, is
\[
                 \ker(\Sigma:G^{18}\to G)\rtimes S_{18},
\]
up to choices of origins in the eighteen torsor fibres.
\end{lemma}
\begin{proof}
The image of $H^0(A,L\otimes\alpha^{-1})$ on $C$ is a very ample
linear system. Its projective dimension is eight, except that it
is seven when the restriction has its one-dimensional kernel
$H^0(A,\alpha^{-2})$. The original coefficient space maps onto
the image system with geometrically integral affine fibres off
the possible centre. It suffices to calculate the monodromy for
the image system on $C$.

Here are the characteristic-two tangency details. Every translate
of $L$ separates second jets on the surface, by
Theorem~\ref{thm:ppav-three-theta-jets}; consequently its restricted
system separates second jets on $C$ at every point. On $C\times C$
off the diagonal, the joint curve first-jet rank is four outside
a set of dimension at most one, is at least three outside a finite
set, and is at least two everywhere. Indeed the corresponding
surface ranks are six, at least five, and at least four. Passing
to curve jets discards only the two normal derivative coordinates.
The rank-six exceptional locus has finite fibres over either
fixed surface point by
Corollary~\ref{cor:all-theta-fixed-point-jets-candidate}, so its
intersection with $C^2$ has dimension at most one. The possible
surface rank-four locus is contained in a finite union of curves
\[
                 (P,Q)=(P,s-P),\qquad P+s\in\Theta,
                 \qquad 3s=0,
\]
after a translation for the chosen twist. Its intersection with
$C^2$ is finite: the first coordinate lies on a theta translate,
which cannot equal the smooth integral curve $C$ of numerical
class $3\Theta$.

Thus the incidence of two distinct tangencies has codimension at
least two in the coefficient space. Its three strata have point
dimensions at most $2,1,0$ and impose at least $4,3,2$ independent
linear conditions, respectively. Contact of order at least three
at one point also has codimension at least two, because three
curve jets are independent everywhere. The single-tangency
incidence is an irreducible projective bundle over $C$; a section
has finitely many zeros, so its map to coefficient space has finite
fibres. Its image is an irreducible divisor, whose geometric generic
member has exactly one double zero and sixteen other simple zeros.
In particular the unlabelled $S_{18}$ assertion is an instance of
the quasireflexive hyperplane theorem
\cite[Theorem~2 and Proposition~3.1]{Entin2021} with its hypothesis
verified here; we also give the labelled argument.

At the generic double zero the degree-two root cluster is
$z^2+bz+c=0$. The independent constant and linear jets give
independent coefficient parameters $b,c$. Along the reduced
discriminant $b=0$ the residual $c$ is not a square, including
after separable closure of the residue field. The equation
$y^2+y=c/b^2$ therefore has nontrivial inertia of order two.
This accounts for the possibly purely inseparable tangent-point
residue extension. The etale torsor is locally trivial on this
cluster, so inertia is a labelled transposition and fixes the
other sixteen torsor fibres pointwise.

The incidence of a point of $\widetilde C$ and a hyperplane
through its image is an irreducible projective bundle over
$\widetilde C$. Its map to coefficient space is finite, since a
nonzero member of the image system cannot vanish identically
on $C$. Hence lifted-point monodromy is transitive. It commutes
with the global deck transformations, and therefore lies in
$G^{18}\rtimes S_{18}$. The Galois closure is branched only over
the discriminant. Quotienting by the normal subgroup generated
by its divisorial inertia gives a cover unramified in codimension
one. Purity \cite{Stacks0BMB} and geometric simple connectedness
of projective space make that quotient trivial, exactly as in
Lemma~\ref{lem:mixed-lifted-intersection-wreath-candidate}.
Thus the image is generated by labelled transpositions.
Lemma~\ref{lem:labelled-transposition-wreath-candidate} gives
the displayed group. Finally the smooth sharp members form a
nonempty dense open by the rank theorem. Restricting to that open
does not change the connected components of any finite etale
cover of the normal coefficient open.
\end{proof}

\begin{lemma}[All boundary gluing columns]
\label{lem:nonordinary-ppav-two-boundary-gluing-candidate}
For fixed $\alpha$, vary independent smooth sharp transverse
$C\in|L\otimes\alpha|$ and $D\in|L\otimes\alpha^{-1}|$.
Their union is in the maximal relative-rank open of $|L^2|$.
Write $A_0=H^1_{\et}(A,\mathbf F_2)$ and $H_C,H_D$ for the two
component cohomologies. The boundary representation has
\[
 \begin{gathered}
 0\longrightarrow T\longrightarrow V\longrightarrow R\longrightarrow0,
 \qquad T=\mathbf F_2^{18}/\mathbf F_2\mathbf1,\\
 R=(H_C\oplus H_D)/A_0,
 \qquad \dim R=16+f,\qquad
 T^0=\{(x_i):\sum x_i=0\}/\mathbf F_2\mathbf1.
 \end{gathered}
\]
It contains an actual $\GL_8\times\GL_8$ acting on two natural
eight-dimensional summands and fixing a complement, and all
translations $\operatorname{Hom}(R,T^0)$. Every nonzero invariant
subspace of this boundary representation contains $T^0$.
In its block frame the span of the two natural summands and $T^0$
is precisely the global common covector kernel $K$.
\end{lemma}
\begin{proof}
Each component has two-rank $8+f$, and the eighteen nodes add
seventeen unit graph directions. Subtracting the ambient rank
$f$ gives $2(8+f)+17-f=33+f$, the maximal relative rank. The
normalization sequence gives the displayed exact sequence.
Restriction from $A$ injects into either component cohomology;
on coherent cohomology this follows from
$H^1(A,(L\otimes\alpha)^{-1})=0$, and hence holds on unit parts.

Put $U_C=H_C/A_0$, $U_D=H_D/A_0$. Their independent relative
images are $\GL_8\times\GL_8$: both variables range over a dense
open of a product, and every Picard twist is a translate of $L$.
If $f=0$ the full component images are already $\GL_8$.
If $f=1$, a component image fixes $A_0$ and surjects onto $\GL_8$;
its unipotent kernel is zero or the entire natural-dual module.
The proved vanishing $H^1(\GL_8,\mathrm{nat}^*)=0$ splits its
complement. Thus in either case the component image is perfect
and has only $\GL_8$ as a nonabelian simple composition factor.
The image on $R$ is likewise a split extension of
$L_0=\GL_8\times\GL_8$ by a submodule of
$(U_C^*\oplus U_D^*)\otimes A_0$; it is perfect. The two
nontrivial simple modules here are nonisomorphic, and the same
first-cohomology vanishing proves the assertion about splitting.

The preceding lemma gives $S_{18}$ on $T$. This group and the
image on $R$ have no common nontrivial quotient: the latter is
perfect with only $\GL_8$ nonabelian simple factors, while
$S_{18}$ has factors $A_{18}$ and $C_2$. For example $31$ divides
$|\GL_8(\mathbf F_2)|$ but not $|S_{18}|$.
Goursat's lemma therefore gives the product joint image.
Take the inverse image of $1\times L_0$. Its translation kernel
lies in $\operatorname{Hom}(R,T)$, which under $L_0$ is a direct
sum of natural-dual and trivial modules. All its quotients have
zero first cohomology: use natural-dual vanishing and perfectness
for the trivial modules. The graph-of-a-cocycle argument gives
an actual block-diagonal $L_0$, fixing $T$ and an ambient
$f$-dimensional complement in $R$.

Let $J\subset\operatorname{Hom}(R,T)$ be the kernel of the action
on $T\oplus R$. Fix $C$ and its maximal connected elementary
two-cover, with group $G=H_C^*$ of dimension $8+f$. Its lifted
intersection image is $G^{17}\rtimes S_{18}$. Its only nonabelian
simple factor is $A_{18}$, so its joint image with the moving
$D$ cohomology is their product by Goursat and perfectness of
the latter. Hence a pure lifted-node translation $(g_i)$ with
$\sum g_i=0$ can be chosen with trivial moving-$D$ cohomology
action. It is an element of $J$, and on the injected $H_C\subset R$
its gluing column is
\[
                  \xi\longmapsto (\xi(g_i))_i
                                    \bmod\mathbf F_2\mathbf1.
\]
Indeed extend the fixed etale two-torsor represented by $\xi$
on $C$ using the trivial torsor on $D$ and chosen node gluing
values. Continuing the choices changes them by these node
evaluations; changing either component trivialization removes
the diagonal vector. This directly identifies the gluing action.

These columns exhaust $\operatorname{Hom}(H_C,T^0)$. Under $L_0$,
$\operatorname{Hom}(R,T)$ is semisimple, with isotypic pieces
for $U_C^*$, $U_D^*$ and, if $f=1$, the trivial ambient column.
Consequently their surjectivity on $H_C$ isolates all maps from
$U_C$ and $A_0$ to $T^0$, with the other columns zero. Reversing
$C,D$ supplies all maps from $U_D$. This proves
$\operatorname{Hom}(R,T^0)\subset J$.

If an invariant subspace has a nonzero projection to $R$, these
translations put all of $T^0$ in it. If it is contained in $T$,
choose a nonconstant representative of a nonzero vector. A
transposition exchanging unequal coordinates gives its difference
$e_i+e_j$, whose $S_{18}$-orbit spans $T^0$.
Finally every global invariant covector kills both natural
blocks and, by the translations, $T^0$. Their dimensions sum
to $8+8+16=32=\dim K$, proving the final assertion.
\end{proof}

\begin{lemma}[Two partitions of a six-theta curve]
\label{lem:nonordinary-ppav-two-six-theta-comparison-candidate}
The geometric representation on $K$ is irreducible. Every
nonzero invariant subspace of $V$ contains $K$.
\end{lemma}
\begin{proof}
Choose general $a_1,\ldots,a_6\in A$ with $\sum a_i=0$.
Their theta translates are distinct, all pairs meet transversely
in two points, and there are no triple intersections. To see
nonemptiness under the sum constraint, pair differences remain
surjective parameters, and the theta difference map is generically
etale. For three chosen indices, vary their three translations
freely and compensate the sum using a fourth index. The incidence
of a common point has dimension at most five in their
six-dimensional translation space, so a generic triple avoids
it. There are only finitely many pairs and triples. The theorem
of the square gives a nodal curve
\[
                    Y=\bigcup_{i=1}^6(\Theta+a_i)\in|L^2|.
\]
Its graph has six vertices and two edges $e_{ij,0},e_{ij,1}$
for every unordered pair $i,j$. Thus it has thirty edges and
graph cohomology of dimension $30-6+1=25$. Its relative stable
rank is $25+6f-f=25+5f$, below the generic rank $33+f$.

For any partition into triples $I\sqcup J=\{1,\ldots,6\}$,
put $C_I=\bigcup_{i\in I}(\Theta+a_i)$ and similarly $D_J$.
They belong to $|L\otimes\alpha_I|$ and
$|L\otimes\alpha_I^{-1}|$. The two twists need not be trivial
and may differ between partitions. The two triple curves are
smooth at their eighteen cross intersection points, so these
points and their graph subgroup $T$ remain finite etale on a
neighbourhood in the complete product of the two coefficient
systems, including at this nodal pair.

Let $0\ne S\subset K$ be globally invariant. Extend it from
the smooth maximal-rank open to the entire coherent maximal-rank
open. This is legitimate because the full constant-rank Lang
scheme is finite etale there and a dense open of a normal
integral base has the same connected components of such a
cover. The full coherent relative Frobenius bundle is regular
on coefficient space: all these ample divisors are connected,
their arithmetic genus is constant, and the negative ample
cohomology vanishing injects ambient $H^1(\mathcal O_A)$.
Take the closure $\overline S$ in its full relative Lang scheme.
Lemma~\ref{lem:normal-base-clopen-Lang-closure-candidate} says
that it is an open-and-closed subgroup, including over $Y$.

On the generic maximal-rank locus of each boundary product,
the preceding lemma puts its entire cross-heart $T^0$ in $S$.
Since $T^0$ is finite etale and extends over the nodal pair,
the clopen closure lemma puts its special fibre in the same
space $\overline S_Y$, independently of the partition. This
uses only those extending finite Lang sections; there is no
identification of the full rank-dropping special fibre with $V$.
The graph space injects into relative cohomology: ambient
$H^1(\mathcal O_A)$ restricts isomorphically to the cohomology
of each principal theta normalization, whereas graph classes
restrict to zero on every normalization. The same holds on
unit parts.

Compare the partitions
\[
        \{1,2,3\}\mid\{4,5,6\},\qquad
        \{1,2,4\}\mid\{3,5,6\}.
\]
The second cross-heart contains
\[
             v_C=e_{13,0}+e_{23,0},\qquad
             v_D=e_{45,0}+e_{46,0}.
\]
Both vectors have even cross-edge sum. On the first partition,
$v_C$ has a nonzero projection to the internal graph of the
first triple and zero projection to the other; $v_D$ has the
opposite property. Nonvanishing can be read on a parallel-edge
two-cycle: the coordinate on $e_{13,0}$ differs from that on
$e_{13,1}$, and likewise for $e_{45,0},e_{45,1}$. In particular
neither internal projection is a vertex coboundary.

Both vectors belong to $\overline S_Y$. Pull back to the strict
henselization of the first boundary parameter space at this
pair. Each point of the etale scheme $\overline S$ has a unique
section with its specified closed value. Hence these two
sections belong to $S$ on the nonempty generic maximal-rank
part of that family. Restriction to each triple curve, followed
by the ambient coherent quotient, is a regular Frobenius map
of vector bundles. Its two nonzero graph projections at $Y$
remain nonzero generically. This proves that $S$ has nonzero
projection to both relative natural systems $U_C,U_D$ of the
first boundary. It does not assert surjectivity of a kernel
under arbitrary rank-dropping specialization.

The two actual independent $\GL_8$ blocks now force the image
of $S$ in $R$ to contain $U_C\oplus U_D$. Its kernel already
contains $T^0$. Thus $\dim S\ge16+16=32=\dim K$, and $S=K$.
For a general nonzero invariant subspace of $V$, the boundary
lemma first supplies $T^0\subset K$ in its intersection with
$K$. The irreducibility just proved then puts all of $K$ in it.
\end{proof}

\begin{theorem}[The full degree-two pointwise stabilizer]
\label{thm:every-nonordinary-ppav-degree-two-full-candidate}
Let $(A,M)$ be a principally polarized abelian surface over any
perfect field of characteristic two, with principal polarization
represented by $M$ over that field and geometric two-rank
$f\in\{0,1\}$. Put $L=M^3$. On the smooth maximal-rank open of
$|L^2|$ the relative unit system has rank $33+f$ and
\[
 G_{\mathrm{geom}}=G_{\mathrm{arith}}
   =\operatorname{Hom}(\mathbf F_2^{\,1+f},\mathbf F_2^{32})
                      \rtimes\GL_{32}(\mathbf F_2).
\]
This is the full pointwise stabilizer of its entire invariant
covector space. Their common kernel is irreducible and is
contained in every nonzero invariant subspace.
\end{theorem}
\begin{proof}
First work over an algebraic closure and suppose the principal
surface is indecomposable. The comparison lemma gives the
irreducibility of $K$ and absence of invariant vectors in $V$:
a fixed line could not contain the sixteen-dimensional $T^0$.
The two actual natural blocks and $T^0$ form all of $K$, and
the boundary contains every rank-one map from either natural
block to $T^0$, zero on the other summands.

Let $N$ be the normal subgroup of the image on $K$ generated
by its transvections. Its restriction is semisimple: the
nonzero $N$-socle is invariant under the whole irreducible
group. Every isotypic multiplicity is one. Indeed a nontrivial
transvection has rank one and therefore can act nontrivially
on an isotypic component only with multiplicity one; a component
fixed by all generators would give a global invariant
$N$-fixed space. The first natural $\GL_8$ block belongs to
one simple $N$-summand, because its natural constituent has
multiplicity one in $K$ restricted to that actual block and
the block lies in $N$. The rank-one maps from it to all vectors
of $T^0$ put $T^0$ in that same summand. The analogous maps
from the second natural block put the second block there as
well. Thus $N$ is irreducible on all of $K$.

The previously used transvection classification
\cite[Theorem~3.6, printed p.~22]{KantorSeress2001} gives
$N=\GL(K)$ or an alternative preserving a nondegenerate
alternating form. The actual natural $\GL_8$ block has no
invariant bilinear form, and its pairing with its fixed
complement is zero since it has no invariant linear functional.
Such a block would lie in the radical of that alternating
form. Therefore the image on $K$ is $\GL_{32}(\mathbf F_2)$.

Put $Q=V/K$, the constant space of dimension $1+f$. The
translation kernel in $\operatorname{Hom}(Q,K)$ is
$K\otimes Q'$ for a subspace $Q'\subset Q^*$. Modulo this
kernel the subgroup is the graph of a cocycle, and the proved
vanishing $H^1(\GL_{32},K)=0$ makes it a linear complement.
If $Q'\ne Q^*$, a nonzero direction of $Q$ annihilated by
$Q'$ would then have a fixed lift in $V$, a contradiction.
Thus the translation kernel is all of $\operatorname{Hom}(Q,K)$.

If the principal surface is geometrically a product, the
already proved all-degree supersingular-product and mixed-product
theorems give the same degree-two formula. These cases exhaust
the two possible nonordinary product types. Finally
Theorem~\ref{thm:nonordinary-ppav-all-covectors} proves that the
entire geometric covector space descends pointwise over every
perfect base field. The arithmetic image is consequently
contained in the displayed stabilizer and contains its full
geometric subgroup, so equality follows.
\end{proof}

\begin{remark}[Scope and exact graph supplement]
\label{rem:nonordinary-ppav-degree-two-scope-candidate}
The six-theta comparison takes place at relative rank $25+5f$,
not at generic rank $33+f$. Its use of clopen Lang closure and
regular component projections is essential. Opposite Picard
twists are allowed separately for the two partitions.
For an ambient hypersurface family the conclusion transfers
whenever restriction to the complete linear system is surjective,
by the previously proved geometrically integral coefficient-fibre
argument. No assertion about all higher degrees on a fixed
nonproduct surface is inferred here. The exact finite supplement
\path{agents/monodromy/code/nonordinary_degree_two_six_theta_graph.py}
checks the doubled six-vertex graph, both cross-hearts and the
two internal projections. It is a finite verification of the
displayed graph calculation, not a substitute for the geometric
or finite-Lang arguments.
\end{remark}
